\documentclass[11pt,a4paper]{article}
\usepackage[margin=22mm]{geometry}
\usepackage{amsmath,amssymb,amsthm,booktabs,microtype}
\usepackage[T1]{fontenc}
\usepackage{lmodern}
\usepackage[colorlinks=true,linkcolor=blue,urlcolor=blue]{hyperref}
\hypersetup{pdftitle={Compatible defects, retained prefixes, and the audited bound 89},pdfauthor={Sum-of-Subset research project}}
\newcommand{\FS}{\operatorname{FS}}
\newcommand{\Z}{\mathbb Z}
\newcommand{\T}{\mathcal T}
\newtheorem{theorem}{Theorem}[section]
\newtheorem{lemma}[theorem]{Lemma}
\newcommand{\RootGroups}{30,396}
\newcommand{\RootSystems}{440,904}
\newcommand{\FinalSystems}{441,243}
\newcommand{\Vectors}{1,534}
\newcommand{\MinGap}{1/748}
\newcommand{\NegativeControls}{50}
\title{Subset Sums in a Fourfold Sumset\\\large Compatible defects, retained prefixes,\\and the audited bound $N\le89$}
\author{Sum-of-Subset research project}
\date{18 October 2026}
\begin{document}
\maketitle
\begin{abstract}
We exclude cardinality 90 and continue the independently audited bound
$N\le90$. Three changes shorten or strengthen the proof. First, retaining
the elementary compatibility $\mu\ge\rho$ lets the existing variable-threshold
hole theorem close the inner lattice comparison directly; no hole theorem
for a deleted prefix is needed at this stage. Second, a budget test rules
out disjoint subset-sum translates in each of the final two prefix steps.
Their non-progression defects and their retained holes give stronger span
and mass inequalities. Third, the 110 fixed covering scores carry 1,510
provable order relations, generated by 180 cover relations. These induce
monotonicity constraints between the loss indicators. A further 46 sign
clauses summarize 7,945 valid pair and triple sign implications. The accompanying
standard-library checker reconstructs every finite case and verifies exact
integer certificates. The archived bound 84 remains outside the completed
independent replay.
\end{abstract}

\section{Result and inherited quantities}
For a finite real set $A$, write $\FS(A)$ for its subset sums, including
zero. The sumset $4A=A+A+A+A$ permits repetitions. Let $N$ be the largest
size of a set satisfying $\FS(A)\subseteq4A$.
\begin{theorem}\label{main}
Every feasible set has at most 89 elements. Consequently $71\le N\le89$.
\end{theorem}
The lower construction is unchanged. The preceding audited result
$N\le90$ \cite{previous} reduces the new task to cardinality 90. Such a
counterexample must contain zero, since otherwise adjoining zero preserves
feasibility and gives 91 elements. This uses the preceding global bound,
not an assertion that arbitrary subsets remain feasible.

We use the corrected integralization, two peeling steps, and necessary
covering reduction in \cite{foundation,previous}. With
$E=P\cup\{0\}\cup(-Q)$ the twice-peeled core, set
\[
 d=82,\quad e=74,\quad a=|P|,\quad b=|Q|,\quad a+b=m=73,
 \qquad I=(a-36)(a-37).
\]
Writing $T_r=r(r+1)/2$, the endpoint constants are
\[
 \Lambda=1+T_a+T_b=1370+I,\quad \Omega=531-I,
 \quad C=2m+4e+14-\Lambda=-914-I.
\]
For a positive $r$-set $X$, put
$\mu_X=|\FS(X)|-1-T_r$ and
$\rho_X=\mu_X+2-2\delta_0(X)$, where $\delta_0$ counts diagonals
absent from the restricted twofold sumset of $X\cup\{0\}$.
We write $\mu=\mu_P+\mu_Q$ and $\rho=\rho_P+\rho_Q$.
The maximum of $X$ is always defective, so
\begin{equation}\label{compatibility}
 \mu_X\ge\rho_X.
\end{equation}
The inherited universal bounds are $\rho_X\ge-8$ and, for $r\ge14$,
$\rho_X\ge0$. Non-progression support has $\mu_X\ge r-3$;
its available defect floor $f(r)$ is $r-14$ for $14\le r<17$ and
$r-5$ for $r\ge17$ \cite{sharp}.

The exact budget, before discarding any nonnegative slack, is
\begin{equation}\label{budget}
 \rho+2\vartheta+c+\T=\Omega,\qquad
 c=2|2D|+18d+85-|\FS(D)|,\qquad
 \T=\kappa+\ell_1+\ell_2+g.
\end{equation}
Here $\kappa=|\FS(P)-\FS(Q)|-|\FS(P)|-|\FS(Q)|+1$.
All displayed slacks except potentially $\rho$ are nonnegative.
The deficient-sign check is $1+T_{74}=2776>24\cdot90+77=2237$.
The paired bootstrap gives odd $q\le51$ from $q^2\le2701$, hence
$a,b\ge11$. Side sizes 11, 12, 13 give respectively
$\Omega=-119,-69,-21<-8$ and are impossible. Thus $a,b\ge14$.

\section{The lattice gate without prefix hole estimates}
We first establish the common core lattice, then propagate it to $D$.
All numerical inequalities in this section are recomputed for all 46
oriented values $14\le a\le59$.

\subsection{Non-progression and progression supports}
If both supports are non-progressions, the sufficient strict inverse test
used in \cite{previous} has smaller-side margin
\[
 T_r-2-\Omega+(r-3)+f(r)+f(73-r)\ge7,
 \qquad r=\min(a,b).
\]
The larger baseline margin is at least 170. The inverse theorem therefore
gives a common lattice, $h_P,h_Q\le\kappa$, and the strict side-hole
conditions $h_X<|\FS(X)|-1$. Normalize this lattice to the integers; then
$h_X=\sum X+1-|\FS(X)|$ counts the missing support points.
The strict side-hole conditions also force each side to be primitive:
a side of coefficient gcd at least two would have
$h_X\ge|\FS(X)|-1$.
Then $z=h_P+h_Q-\kappa$ satisfies $0\le z\le\min(h_P,h_Q)$.

For progression supports we repeat the elementary gap repair of
\cite{previous} with the new constants. For $W=\{w_1<\cdots<w_v\}\subset\Z$ and integer $u\ge1$,
\[
 |W+[0,u-1]_{\Z}|=v+u-1+
 \sum_{i=1}^{v-1}\min(w_{i+1}-w_i-1,u-1).
\]
In $s$ residue classes the excess is at least $(s-1)(u-1)$.
A progression side of size at least 37 has $u-1\ge703>\Omega$,
forcing one residue and no long gap. It gives $z=0$.
For a smaller progression side $a\le36$, the verified bounds are
\[
 \Omega<2T_a,\qquad \Omega-T_a\le75.
\]
A unique long support gap forces the largest coefficient to exceed the
sum of its predecessors, by complement symmetry of the unique gap.
A sole off-lattice coefficient likewise gives two disjoint support copies.
Either would force $\rho_Q\ge T_{b-2}+2-2b\ge558>531$.
Thus only the two-residue exception needs attention.

In that exception $\kappa\ge u-1\ge T_a$, so $a\ge18$ and
$u\ge T_a+1>2b$. At least two coefficients occupy the nonzero
residue, whose order is two. Normalize the half-step and write
$U=2[0,u-1]$, with $w$ the least odd coefficient of $Q$.
The same two-fiber identities as in \cite{previous} give
\[
 \kappa=u-1+h_0+h_1,\quad h_P=u-1,\quad
 h_Q=\kappa-u+w,\quad z=w-1.
\]
No fiber gap reaches $u$. If $s\ge2$ coefficients are odd, distinctness
gives $\sum Q\ge(b-s)(b-s+1)+sw+s(s-1)$.
For $w>u$, comparing this with
$\mu_Q+\kappa\le\Omega+2b-2$ has positive margin at least 301,
so $w\le u$ and both side holes are at most $\kappa$.
For $w\ge15$, the lower bound
$\sum Q-w-T_b\ge7b-289/8$ exceeds the allowed $2b+72$.
Therefore $w\le13$ and $z\le12$.
The pre-lattice span comparison becomes
\[
 |\FS(D_1)|-2|\FS(E)|
 \le2513+2a-4T_a-2T_b\le-1089.
\]
The outer margin is at least 521. Both peeled layers are thus integral
on the half-step lattice before the sharp peeling estimate is used.
Now $\T\ge\kappa+154-3z\ge T_a+118$ contradicts
$\Omega-T_a\le75$, with margin at least 43.

The surviving progression cases have $z=0$. The general span estimate
from \cite{previous} gives the required inner comparison $C+\Omega<0$.
For either support type, the boundary $a=14,b=59$ uses the universal
prefix floor $-8$: the span constant increases from 2 to 10 on that side. The inner comparison is still at most
$C+3\Omega+16=-1329$. After lattice propagation,
$\T\ge52>\Omega=25$ excludes this boundary.
In every surviving $z=0$ progression case,
$\kappa\le\Omega-154$ and $T_r-\kappa\ge79$ supply the strict
side-hole bounds as well.

\subsection{Keeping support and defect estimates compatible}
For the non-progression case, the core lattice and strict side-hole
bounds are already known. The two inherited inner comparisons are
\begin{equation}\label{inner}
 |\FS(D_1)|-2|\FS(E)|\le
 \min\{C+\Omega+\rho-\mu+2z,\ C+3\Omega-\rho-\mu\}.
\end{equation}
By \eqref{compatibility}, the first is at most $-383-2I+2z$, hence
strict for $z\le191+I$. In the remaining range $z\ge192+I$, the
basic non-progression floors already settle every orientation except
$25\le a\le48$.

Let $\mathcal R(r,z)$ be the previously proved variable-threshold
persistent-hole bound \cite{previous}, checked by
\path{verification/nonap_hole_rho_schedule.py}. Equation
\eqref{compatibility} gives the direct bound
\begin{equation}\label{directgate}
 C+3\Omega-\rho-\mu
 \le C+3\Omega-2\bigl(\mathcal R(a,192+I)+
                            \mathcal R(b,192+I)\bigr)<0.
\end{equation}
Its maximum is $-75$. At $(a,b)=(36,37)$ the two defect floors are
184 and 193, so the right side is $679-2(184+193)=-75$.
This uses neither an effective-hole estimate for a deleted prefix nor
prefix integrality beyond the already established core lattice.
The outer margin is $\Lambda-(4d-10+\Omega)=521+2I>0$.
Both peeled layers are integral. Sharp peeling now supplies
\begin{equation}\label{peeling}
 w\le4z,\quad \T\ge\kappa+154-3z,
 \qquad\T\ge\max(52,z,154-2z).
\end{equation}

The size-15 side has no surviving scalar state. For size 16 the coarse
positive-hole candidates have $49\le z\le57$; applying
$\mathcal R(57,z)$ only to the other side excludes all 18 oriented
candidates. No size-16 call to a lemma requiring size 17 is made.
All surviving positive-hole states have $a,b\ge19$.

\section{Retaining the last two prefixes}
\begin{lemma}\label{retention}
Let $X=\{x_1<\cdots<x_r\}$ be a primitive positive integer set with
non-progression subset-sum support. Let $q\ge1$ be an integer with
$r-q\ge17$, and suppose
\[
 \rho_X\le B<T_{r-q}-r-1.
\]
For every $1\le j\le q$, the prefix $X_{r-j}$ is primitive and
has non-progression support. Moreover
\[
 h_{X_{r-j}}\ge 2^{-j}h_X,\qquad
 \rho(X_{r-j})\ge r-j-5,\qquad
 \mu(X_{r-j})\ge r-j-3.
\]
\end{lemma}
\begin{proof}
Write $y_i=|\FS(X_i)|-|\FS(X_{i-1})|-i$.
The inherited recurrence is $\rho_i\ge\rho_{i-1}+y_i-2$;
non-progression support at size $i\ge6$ implies $y_i\ge1$.
If the last translates at size $i$ are disjoint, then
$y_i=|\FS(X_{i-1})|-i\ge1+T_{i-1}-i$.
Since $\rho_{i-1}\ge0$, this yields
$\rho_i\ge T_{i-1}-i-1$.

Proceed backwards. At a retained non-progression prefix of size $i$,
the already traversed steps give $\rho_r\ge\rho_i-(r-i)$.
Disjointness at the next backward step would therefore imply
$\rho_r\ge T_{i-1}-r-1\ge T_{r-q}-r-1>B$.
Thus the translates intersect. Their intersection forces the last
coefficient onto the prefix coefficient lattice, so their gcds agree.
If the preceding support were an arithmetic progression, its intersecting
translate would have the same step and their union would be a progression.
Both primitivity and non-progression therefore persist backwards.

The two support hulls overlap. Every hole of the union belongs to a
hole of at least one of the constituent supports, so $h_i\le2h_{i-1}$.
Induction gives the claimed hole bound, and the inherited sharp
non-progression and support bounds give the other two conclusions.
\end{proof}

For each surviving positive-hole state, set
$t(z)=\max(52,z,154-2z)$ and use the established lower defect
of the other side. The budget bounds
$\rho_X\le\Omega-\rho_{\rm other}-t(z)$.
The exact minimum surplus in Lemma~\ref{retention} is 68 for $q=1$
and 44 for $q=2$. All surviving sides have size at least 19, so the
sharp prefix floors apply. These inequalities are checked for each
state, including low holes; they are not inferred from sampled sets.

Let $Y=X\setminus\{R\}$, where $R=\max X$.
The exact hole identity and the defect recurrence give
\begin{align}
 R&=r+(\mu_X-\mu_Y)+h_X-h_Y,\\
 R&\le r+\rho_X+h_X+2-\rho_Y-h_Y.\label{spanidentity}
\end{align}
Applying Lemma~\ref{retention} on both sides yields
$\rho_{P'}+\rho_{Q'}\ge m-12=61$ and
$h_{P'}+h_{Q'}\ge z$. Combining \eqref{spanidentity} with the
budget and peeling gives, for $j$ empty ordinary units and
$\delta=c-82j$,
\begin{equation}\label{newspan}
 R_p+R_q+\delta-3z\le393-I-82j.
\end{equation}
The preceding unretained span was
$R_p+R_q+\delta-4z\le454-I-82j$.
Thus the new row gains $61+z$ at the same actual state.
It is used only when $z>0$.

There are also new mass rows. On one side let
$S$ be the total positive mass of $D$, and
$(R,x_1,x_2,x_3,M)$ its core maximum and four successive outer
magnitudes; let $C_1$ be the second largest core element.
The one-prefix and two-prefix bounds respectively imply
\begin{align}
 2M+8R-2S+z&\le-2T_{r-1}-12-2(r-4),\\
 2(R+x_1+x_2+x_3+M)-2S+z&\le-2T_{r-1}-2(r-4),\\
 4(R+C_1+x_1+x_2+x_3+M)-4S+z
   &\le-4T_{r-2}-4(r-5).\label{secondmass}
\end{align}
For example, the remaining prefix mass is
$T_{r-1}+\mu_Y+h_Y\ge T_{r-1}+r-4+z/2$;
the next prefix uses $h\ge z/4$. These explain each factor 2 or 4.
All masses and coefficients in the checker remain exact.

The helper \path{verification/retained_prefix.py} validates the budget
witnesses and additionally checks 2,916 actual finite sets, including
1,333 intersecting and 1,583 disjoint steps, and 492 retained-prefix
instances. The factor-two hole loss is sharp: the prefix
$\{1,3,4,\ldots,19\}$ has two holes, while adjoining 187 gives four
holes with intersecting translates. These finite checks supplement the
general proof; they do not establish it by enumeration.

\section{Fixed scores as an ordered family}
The inherited zero-hole refinement charges four high fixed labels when
the largest own center is empty and a lower own center contributes.
Only the completely integral four-empty mode
\texttt{pair1,pair2,p0,q0} uses the fixed-score extension here.
Both sides are charged, so eight labels are already removed.
For each other fixed score $f_v$, let $H_v$ indicate $f_v>S_v$.
For all 110 labels, let $N_v$ indicate $f_v\le0$.
The 26 disjoint deletion groups require at least one nonpositive score
per group. The additional fixed loss is therefore at least
\[
 \sum_{v\text{ not charged}}H_v+\sum_vN_v-26.
\]
The high and nonpositive events are disjoint; the eight separately
charged labels are omitted from the high sum. This is the same valid
deficit allocation as in \cite{previous}.

The new observation is that these indicators are ordered.
Write every score in the eight positive outer-gap coordinates.
If the coefficient vector of $f_v-f_u$ is componentwise nonnegative,
then $f_u\le f_v$ in every actual configuration. Consequently
\begin{equation}\label{scoreorder}
 N_v\le N_u,\qquad
 H_u\le H_v\quad\hbox{when $u,v$ have the same target mass}.
\end{equation}
Also $H_v+N_v\le1$. These statements concern actual indicators,
so adding them as necessary linear rows is valid even though the
relaxation permits fractional indicator values.

The checker regenerates every score from the original covering labels,
then forms and checks every coefficient difference. There are 1,510
ordered pairs, all with the same respective target side, and 180 cover
relations. It verifies that the transitive closure of those 180 relations
equals the complete computed order. Only cover relations need new rows.
This compression is a proof about the coefficient cone; it does not rely
on checking a finite sample of outer elements.

There is a second, complementary use of the order. Whenever two or three
score forms have componentwise nonpositive total gap coefficients,
at least one of their actual scores is nonpositive. Thus the associated
indicators satisfy $\sum_{v\in J}N_v\ge1$.
The checker reconstructs 209 such pairs and 7,736 such triples.
A clause may be lowered in the score order: if a nonpositive score has
a smaller score in the retained clause, that smaller score is nonpositive
too. Remove comparable redundant labels and associate each clause with
its upward closure in the score poset. Of 7,489 distinct upward sets,
46 are inclusion-minimal; their Boolean sign implications imply all
7,945 original pair/triple implications. Each retained clause keeps an
explicit original pair/triple witness and its nonpositive coefficient
sum. The checker verifies both that witness and every order implication.
It adds the 46 necessary linear indicator rows, while retaining the
original 26 disjoint-group rows and their separate counting justification.
The compression claim concerns actual Boolean indicators; no unsupported
linear implication between fractional relaxations is assumed.
The original score
dichotomies remain exhaustive because the relevant coordinates and
target masses have already been proved integral.

\section{Complete covering domain and exact verification}
We retain the eighteen ordinary units, contributor orientations, and all
integer-center witnesses from \cite{previous}. In particular, all
10,880 four-empty modes remain present; the 30 incomplete integer bases
use only the gaps actually proved integral.
The endpoint constants in every model are rebuilt:
\begin{center}\small
\begin{tabular}{lr}\toprule
Quantity&Value\\\midrule
Ordinary size; core size&82; 74\\
Global quantiles; core quantiles&82/81; 78/77\\
Basic, nested, second-core leading thresholds&163, 155, 153\\
Core loss threshold&78\\
Zero-hole span right side&$454-I-82j$\\
Positive-hole retained span right side&$393-I-82j$\\
Cardinality row right side&$1560-82j$\\\bottomrule
\end{tabular}
\end{center}
The exact low-hole envelopes are
\begin{center}\small
\begin{tabular}{lrrrr}\toprule
Family&$z$&Combined defect floor&Support floor&$c+I\le$\\\midrule
Z&0&0&0&377\\
P&1--2&63&$r-3$&318\\
Q&3--4&75&$r+6$&310\\
R&5--14&96&$2r-5$&309\\\bottomrule
\end{tabular}
\end{center}
Each group's integer range is reconstructed from
$I\le\text{cap}-82j$. The positive families transport $z$ to their
own upper endpoint and add the same amount to both mass coordinates.
The new prefix rows preserve this implication: their total coefficient
under this transport is nonpositive. Enlarging a necessary model's
parameter domain to include scalar-impossible cases is harmless.

For high holes, the same variable $z$ is retained in each merged group.
With certified defect floors $r_p,r_q$, put
$E_*=\Omega-r_p-r_q$ and use all four coupled caps
\[
 \delta\le c_{\max}-82j,\quad \delta+z\le E_*-82j,
 \quad\delta\le E_*-52-82j,
 \quad\delta-2z\le E_*-154-82j.
\]
The support floors are now also replaced by
$\max(\hbox{effective-hole support floor},\hbox{defect floor})$,
using \eqref{compatibility}. This reduces the all-active unresolved
high-hole states from 574 to 333 in the same size-90 scalar relaxation.
It is a comparison of necessary relaxations, not itself an endpoint
exclusion. All 1,866 surviving high-hole states have $c\le162<2d$.
The residual cap of each actual group decides which quantile rows apply;
$162$ is not below the deeper leading threshold 153, and $162-82=80$
is not below the core threshold 78.

Full-translate common deficits, trimmed two-sum bounds, all inherited
geometric splits, the score indicators, \eqref{scoreorder}, and the 46 sign clauses complete
the necessary models. Every source-to-group implication is checked
row by row for both coarse and nested models.

The certificate bundle contains \Vectors{} distinct nonnegative integer
vectors covering \RootGroups{} groups, \RootSystems{} integer roots,
and \FinalSystems{} final leaf systems. The minimum normalized strict
gap is $\MinGap$. The checker reconstructs every root, parameter interval,
and both children of every split. For necessary rows $Ax\le b$ with
$x\ge0$, each accepted multiplier vector satisfies
$w\ge0$, $wA\ge0$, and $wb<0$ exactly.
Symbolic quadratic right sides are checked at the true integer extrema;
the exhaustive mode also rebuilds each numerical matrix and compares
every row before checking its certificate. There are \NegativeControls{}
negative controls. The completed numerical cross-check is partitioned into
eight independent batches by \path{experiments/verify_upper89_sharded.py}.
Every reconstructed root occurs exactly once; every batch checks all its
integer matrices, interval certificates, and both children of every tree.
Their merged per-root results and statistics must equal the standard run.
One earlier uninterrupted attempt ended without a completion receipt and
is recorded as incomplete, not as a pass.

From the repository root:
\begin{verbatim}
python -I -S -B verification/upper89.py --output .work/upper89.json
python -I -S -B verification/upper89.py --exhaustive \
  --output .work/upper89-exhaustive.json
python experiments/generate_upper89_certificates.py \
  --output .work/upper89-regenerated.json
\end{verbatim}
The discovery search minimizes the residual loss when fixed-score indicators
are present. At a remaining fractional solution it tries up to eight
violated, already proved dichotomies and prefers immediate closures,
then the strongest child loss bound. Both children are always retained.
These floating-point choices only guide discovery; acceptance of every
leaf requires the exact integer certificate check above.
Only discovery uses NumPy/SciPy. The verifier uses the standard library.
The Windows entry point is \path{verify_upper89_windows.bat}.
Reports, source hashes, compressed full output, independent regeneration,
optimized-Python comparison, exhaustive cross-checks, and the unchanged
upper-90 regression are recorded in
\path{results/2026-10-18/upper89/}.
This is an exact finite certificate proof coupled to the general lemmas
above and the audited foundation. It is not a proof-assistant
formalization or a completed replay of the archived bound 84.

\begingroup\small
\begin{thebibliography}{9}\setlength{\itemsep}{1pt}
\bibitem{previous} Sum-of-Subset project,
\emph{Variable thresholds for persistent holes and the audited bound
$N\le90$}, \path{reports/upper90_from91_20261018.tex}.
\bibitem{foundation} Sum-of-Subset project,
\emph{From the audited bound 115 to 103},
\path{reports/upper103_from115_20261014.tex}.
\bibitem{sharp} Sum-of-Subset project,
\emph{Sharp non-progression defects and the audited bound 94},
\path{reports/upper94_from95_20261016.tex}.
\end{thebibliography}
\endgroup
\end{document}
